\section[untilt、$\mathcal{O}(1)$ 与 Lubin--Tate 理论]{untilt、$\mathcal{O}(1)$ 与 Lubin--Tate 理论（11 月 21 日）}

\srcpage{32}

\subsection{Untilt}

我们固定通常的记号：$E$ 是非阿基米德局部域，$\mathcal O_E \ni \pi$，$\Fq \subset \ol{\FF}_q$，而
$\breveE = W_{\mathcal O_E}(\ol{\FF}_q)[1/\pi]$ 是 $E$ 的最大无分歧扩张的完备化。

对 $S \in \Perf_{\Fq}$ 是完满空间，我们定义了 $Y_{S,E} \twoheadrightarrow X_{S,E} = Y_{S,E}/\phi_S^{\mathbb Z}$。在菱形层面，我们有
\[
  Y_{S,E}^{\diamond} = S \times (\Spa E)^{\diamond}.
\]

我们想陈述「untilt 等于 $Y_{S,E}$ 上次数为 1 的 Cartier 除子」这一命题。

局部地，任何 $S$ 都形如 $\Spa(R,R^+)$，而一个 untilt 将形如 $S^{\#} = \Spa(R^{\#},R^{\#+})$。存在典范满射
\[
  \theta: W_{\mathcal O_E}(R^+) \twoheadrightarrow R^{\#+}.
\]
（这一点此前已经出现过——由完美环的刚性，这样的映射与 $R^+ \to R^{\#+}/\pi$ 相伴。任何这样的映射都唯一地穿过 $R^{\#+}/\pi$ 的逆向极限完美化，而后者带有与 $R^+$ 的等同。）由整完满环（等价地，「完美棱镜」）的一般结构结果，我们将有 $\ker(\theta) = (\xi)$，其中 $\xi$ 是 $W_{\mathcal O_E}(R^+)$ 中的非零因子。

正如上次那样，这诱导出 $S^{\#} \hookrightarrow \Spa W_{\mathcal O_E}(R^+) \setminus \{\pi = 0,\ [\varpi] = 0\}$，其中 $S^{\#}$ 实现为 Cartier 除子 $V(\xi) \subset Y_{S,E}$。

\subsection{闭 Cartier 除子}

\begin{definition}
设 $X$ 是一致解析（uniform analytic）adic 空间\footnote{解析（analytic）指 $X$ 局部地形如 $\Spa(R,R^+)$，其中 $R$ 是 Tate 环。一致（uniform）指 $R$ 上的谱范数是一个范数，等价地 $R^0 \subset R$ 有界。于是 $R^+ \subset R$ 有界。}（例如 $X$ 完满）。$X$ 上的一个\emph{闭 Cartier 除子}是理想层 $I \subset \mathcal O_X$，它局部自由且秩为 1，并且对所有的 affinoid $U \subset X$，映射 $I(U) \to \mathcal O_X(U)$ 的像都是闭的。
\end{definition}

换句话说，局部地 $I$ 由一个元素生成，该元素是非零因子，并且生成一个\emph{闭}理想。这个闭性条件对于让商行为良好是必要的：此时 $(V(I),\mathcal O_X/I)$ 将定义一个 adic 空间。

\begin{warning}
该条件不能在 affinoid 的开覆盖上检验——必须对\emph{每个} affinoid $U \subset X$ 真正地检验。所以要检验 §7.1 中出现的 $(\xi)$ 是闭 Cartier 除子，还需要再做一点工作。
\end{warning}

\begin{proposition}
(1) $V(\xi) = S^{\#} \hookrightarrow Y_{S,E}$ 是闭 Cartier 除子。

(2) 复合映射 $S^{\#} \hookrightarrow Y_{S,E} \to X_{S,E}$ 是闭 Cartier 除子。
\end{proposition}

\begin{definition}[「次数 1 的 Cartier 除子的模空间」]
设 $\Div^1_Y, \Div^1_X$ 是函子 $\Perf_{\Fq} \to \mathbf{Sets}$，它把 $S$ 映到 $Y_{S,E}$（分别地 $X_{S,E}$）上闭 Cartier 除子的集合，这些除子在 $S$ 上局部地形如 $S^{\#} \hookrightarrow Y_{S,E}$（分别地 $S^{\#} \hookrightarrow X_{S,E}$），其中 $S^{\#}/E$ 是 untilt。
\end{definition}

\begin{remark}
这个函子本可以在 $\Perf_{\Fq}$ 上定义。然而后面我们会看到，出于各种理由，我们想要几何地工作。
\end{remark}

\begin{proposition}
我们有如下菱形的等同。

(1) $\Div^1_Y = (\Spa \breveE)^{\diamond}$。

(2) $\Div^1_X = \Div^1_Y/\phi_E^{\mathbb Z} = (\Spa \breveE)^{\diamond}/\phi_E^{\mathbb Z}$。
\end{proposition}

为解释 (2) 中出现的 $\phi_E$，注意 $(\Spa \breveE)^{\diamond} = (\Spa E)^{\diamond} \times_{\Fq} \ol{\FF}_q$。而 $(\Spa E)^{\diamond}$ 是 $\Perf_{\Fq}$ 上 $S$ 的函子，因此带有来自测试范畴上 Frobenius 的 Frobenius。同理，\emph{任何}菱形都有 Frobenius。

\srcpage{33}

\begin{proof}
(1) 按定义，$(\Spa \breveE)^{\diamond} \twoheadrightarrow \Div^1_Y$ 是满射，因为由 $\Div^1_Y$ 参数化的任何 Cartier 除子局部地来自某个 untilt $S^{\#}/E$。由于我们考虑的是 $\ol{\FF}_p$ 上的对象，$S^{\#}$ 上的 $E$-结构自动升级为 $\breveE$-结构。

但反过来，$Y_{S,E}$ 上的闭 Cartier 除子决定 $Z \subset Y_{S,E}$，它在 $S$ 上局部地是 $S$ 的一个 untilt。这些局部 untilt 粘成 $S$ 的整体 untilt $S^{\#}$。从而 $(\Spa \breveE)^{\diamond} \xrightarrow{\sim} \Div^1_Y$。

(2) 在 (1) 中取关于 Frobenius 的商。映射仍然是满射；我们只需核对哪些对象被等同，这留作练习。
\end{proof}

\begin{warning}
注意 $X^{\diamond}_{S,E} = (S/\phi_S^{\mathbb Z}) \times (\Spa E)^{\diamond}$，但 $\Div^1_X = (\Spa \breveE)^{\diamond}/\phi_E^{\mathbb Z}$。这两个商是对不同的 Frobenius 取的。Fargues 把 $\Div^1_X$ 称为「镜像曲线」（mirror curve）。它「只是」一个菱形（即不来自 adic 空间）。

注意 $\Div^1_X$ 的拓扑空间是一个点。它既不是拟分离的，也不是局部空间的（因为它没有有意思的子集，而空间性包含拟分离性）。

口号：「曲线上点的模空间并不等同于曲线本身」。
\end{warning}

\subsection{$\mathcal O(1)$ 与 untilt}

回忆 $\mathcal O_{X_{S,E}}(1)$ 是 $X_{S,E}$ 上对应于等晶体 $(\breveE,\pi^{-1}\sigma)$ 的线丛。特别地，由下降我们有
\[
  H^0(X_{S,E},\mathcal O_{X_{S,E}}(1)) = H^0(Y,\mathcal O_{Y_{S,E}})^{\phi_S=\pi}.
\]

目标是看出：若 $S^{\#}$ 是 $S$ 的任一 untilt$/E$，则 $I_{S^{\#}} \subset \mathcal O_{X_{S,E}}$ 在（对 $S$ 作 pro-étale 局部化之后）同构于 $\mathcal O_{X_{S,E}}(-1)$。\footnote{这些同构一般不会粘合。}

\begin{remark}
在这个意义下，$S^{\#} \hookrightarrow X_{S,E}$ 是一个「次数 1」的除子。后面我们将证明：当 $S$ 是几何点时，$\Pic(X_{S,E}) \cong \mathbb Z$，而这将给出次数的精确概念。
\end{remark}

为证明目标，我们想构造映射 $\mathcal O(-1) \xrightarrow{\sim} I_{S^{\#}} \hookrightarrow \mathcal O$。作一次扭转，这等价于构造映射 $\mathcal O \to \mathcal O(1)$，即构造沿 $S^{\#}$ 消失的元素 $H^0(X_{S,E},\mathcal O(1))$。我们将用 Lubin–Tate 形式群给出 $H^0(X_{S,E},\mathcal O(1))$ 的一个「公式」。

\subsection{Lubin–Tate 理论}

回忆 Lubin–Tate 形式群是一个一维形式群 $G/\mathcal O_{\breveE}$，带有 $\mathcal O_E$ 的作用，使得在 $\Lie(G)$ 上诱导的两个作用一致。

一维形式群律有 $[\pi]_G(X) = \pi X + O(X^2)$，而 $[\pi]_G(X)$（模 $\pi$）的第一个非零项是 $a_{q^h}X^{q^h}$，其中 $h = 1,2,\dots,\infty$。这个 $h$ 称为 $G$ 的\emph{高度}。

\begin{theorem}[Lubin–Tate]
在同构（非常不唯一）的意义下，对任意固定的高度 $h$，Lubin–Tate 形式群是唯一的。
\end{theorem}

\begin{example}
我们关心高度 1 的 Lubin–Tate 形式群。设 $E = \Qp$。此时可以取 $G = \hat{\mathbb G}_m$ 为形式乘法群 $\Spf \breve{\mathbf Z}_p[[X]]$，而 $X +_G Y = (1+X)(1+Y)-1$。
\end{example}

对任何 Lubin–Tate 形式群 $G$，在一般纤维上有到加法群的同构
\[
  \log_G: G \times_{\mathcal O_{\breveE}} \breveE \xrightarrow{\sim} \hat{\mathbb G}_{a,\breveE}.
\]
（特征 0 的所有一维形式群都通过对数同构于形式加法群。）于是 $G$ 上的群律可以写成
\[
  X +_G Y = \exp_G(\log_G(X) + \log_G(Y)).
\]

\srcpage{34}

\begin{remark}
我们可以选取同构 $G \cong \Spf \mathcal O_{\breveE}[[X]]$，使得其对数形如
\[
  \log_G(X) = X + \frac{1}{\pi}X^q + \frac{1}{\pi^2}X^{q^2} + \cdots.
\]
这给出了构造 $G$ 的方法：通过对数把 $\mathbb G_a$ 上的加法搬运到 $G$ 上，并检验系数是整的。
\end{remark}

\subsection{与局部类域论的联系}

对任意 $n \ge 1$，令 $G[\pi^n] \subset G$ 是 $G$ 上乘以 $\pi^n$ 这一映射的核。于是
\[
  G[\pi^n] \cong \Spf \mathcal O_{\breveE}[[X]]/([\pi^n]_G(X)).
\]

\begin{example}
对 $E = \Qp$，如前所述 $G \cong \hat{\mathbb G}_m$，所以 $G[\pi^n] \cong \mu_{p^n}$。

我们有 $G[\pi^n] \times_{\mathcal O_{\breveE}} \breveE = \bigsqcup_{i=0}^{n} \Spec \breveE_i$，其中 $\breveE_0 = \breveE$，且 $\breveE_0 \subset \breveE_1 \subset \cdots \subset \breveE_n \subset \cdots$ 是一列扩张，$\breveE_n/\breveE_0$ 由添加一个本原的 $\pi^n$-挠点得到。
\end{example}

\begin{theorem}
$\breveE$ 的、在 $E$ 上为 abel 的极大扩张是 $\bigcup_n \breveE_n$，并且有典范同构
\[
  \Gal(\breveE_n/\breveE) \cong (\mathcal O_E/\pi^n)^\times.
\]
又，$\breveE_\infty := \bigcup_n \breveE_n$ 的完备化是完满域。
\end{theorem}

\subsection{$G$ 的万有覆盖}

\begin{definition}
我们定义 $\tilde G := \varprojlim_{[\pi]_G} G$（形式概形范畴中的极限）。
\end{definition}

\begin{proposition}
存在同构 $\tilde G \cong \Spf \mathcal O_{\breveE}[[\tilde X^{1/p^\infty}]]$。
\end{proposition}

\begin{proof}
因为 $\tilde G$ 是 $p$-进完备的，并且在 $\mathcal O_{\breveE}$ 上平坦，由这类对象的刚性，只需证明
\[
  \tilde G \times_{\mathcal O_{\breveE}} \ol{\FF}_q \cong \Spf \ol{\FF}_q[[\tilde X^{1/p^\infty}]].
\]
但模 $\pi$ 时，形式群律被安排成 $X \mapsto X^q$，所以这是清楚的。
\end{proof}

我们有投影映射 $f_n: \tilde G \to G$，在环的层面上由 $\mathcal O_{\breveE}[[\tilde X^{1/p^\infty}]] \leftarrow \mathcal O_{\breveE}[[X_n]]$ 给出。可以验证
\[
  \tilde X = \lim_{n\to\infty} X_n^{q^n}.
\]

在一般纤维上，我们有映射 $\tilde G^{\mathrm{ad}} \times_{\mathcal O_{\breveE}} \breveE \to G^{\mathrm{ad}} \times_{\mathcal O_{\breveE}} \breveE \xrightarrow{\log_G} \mathbb G_a$。在 adic 空间的层面上，$G^{\mathrm{ad}} \times_{\mathcal O_{\breveE}} \breveE$ 是开单位圆盘，而 $\tilde G^{\mathrm{ad}} \times_{\mathcal O_{\breveE}} \breveE$ 是无穷覆盖。复合函数由下式给出（练习）：
\[
  \sum_{i \in \mathbb Z} \pi^i \tilde X^{q^{-i}} \in \mathcal O\bigl(\tilde G^{\mathrm{ad}} \times_{\mathcal O_{\breveE}} \breveE\bigr).
\]

\subsection{与整体截面的关系}

设 $S = \Spa(R,R^+) \in \Perf_{\Fq}$。我们有 Fargues–Fontaine 曲线 $X_{S,E}$ 以及线丛 $\mathcal O_{X_{S,E}}(1)$。设 $S^{\#} = \Spa(R^{\#},R^{\#+})$ 是 $S$ 的一个 untilt。

则对数给出 $\tilde G(S^{\#}) = \tilde G(R^{\#+}) \xrightarrow{\sim} R^{\circ\circ}$，它带有到 $H^0(Y_{S,E},\mathcal O_{Y_{S,E}})$ 的映射，把 $X$ 映到 $\sum_{i\in\mathbb Z} \pi^i[X^{q^{-i}}]$。

\begin{proposition}
这诱导出同构
\[
  \tilde G^{\mathrm{ad}}(S^{\#}) \xrightarrow{\sim} H^0(X_{S,E},\mathcal O_{X_{S,E}}(1)) = H^0(Y_{S,E},\mathcal O_{Y_{S,E}})^{\phi_S=\pi}.
\]
在这个同构下，求值映射 $H^0(Y_{S,E},\mathcal O_{Y_{S,E}}) \to R^{\#}$（在 $S^{\#} \subseteq Y_{S,E}$ 处）对应于对数映射
\[
  \tilde G^{\mathrm{ad}}(S^{\#}) = \tilde G(R^{\#+}) \xrightarrow{\log_{\tilde G}} R^{\#}.
\]
\end{proposition}

\srcpage{35}

特别地，
\[
  H^0(X_{S,E},\mathcal O_{X_{S,E}}(1)) \cong \Hom\bigl(S,\Spa \ol{\FF}_q[[\tilde X^{1/p^\infty}]]\bigr).
\]

\begin{remark}
若 $n \le [E:\Qp]$（分别地，若 $E$ 的特征为 $p$ 则对所有的 $n$），则我们可以计算
\[
  H^0(X_{S,E},\mathcal O_{X_{S,E}}(n)) \cong \Hom\bigl(S,\Spa \ol{\FF}_q[[X_1^{1/p^\infty},\dots,X_n^{1/p^\infty}]]\bigr).
\]
口号：「函子 $S \mapsto H^0(X_{S,E},\mathcal O_{X_{S,E}}(n))$ 由 $\Spa \ol{\FF}_q[[X_1^{1/p^\infty},\dots,X_n^{1/p^\infty}]]$ 表示\footnote{「表示」对象并不完全是完满空间。}，后者是一个 $n$ 维完满开单位圆盘。」

然而，若 $n > [E:\Qp]$，则这个函子不可表示。这些函子称为 Banach–Colmez 空间，它们将是菱形的有趣例子。
\end{remark}

\begin{proof}
与 $\log_{\tilde G}$ 的交换性来自对公式的检查。到 $H^0(X_{S,E},\mathcal O_{X_{S,E}}(1)) = H^0(Y_{S,E},\mathcal O)^{\phi=\pi}$ 的映射是显然的：检查公式即可看出 Frobenius 确实乘以 $\pi$。

等特征情形来自直接计算。若 $E = \Fq((t))$，则 $Y_{S,E}$ 是去心开单位圆盘，所以 $H^0(Y_{S,E},\mathcal O)$ 是某些显式的幂级数环，由形如 $\sum_{n\in\mathbb Z} t^n r_n$（其中 $r_n \in R$）的元素组成，并且要求在去心开单位圆盘上收敛。$\phi_S$ 等于乘以 $\pi = t$ 这一条件，等价于 $r_{n+1}^q = r_n$。所以一切都由 $r_0 = r \in R$ 决定，而它受制于条件：$\sum_{n\in\mathbb Z}\pi^n r^{1/q^n}$ 收敛。容易检验这当且仅当 $r \in R^{\circ\circ}$ 是拓扑幂零的。

若 $E/\Qp$，则这可以从 [SW] 中 Dieudonné 理论的一个结果推出。
\end{proof}

\begin{warning}
若 $E$ 是 $p$ 进的，则人们\emph{不能}把 $H^0(Y_{S,E},\mathcal O_{Y_{S,E}})$ 描述为某些和
\[
  \sum_{n\in\mathbb Z} \pi^n[r_n], \quad r_n \in R.
\]
更确切地说，已知对 $H^0(Y_{S,E},\mathcal O_{Y_{S,E}})$ 的所有元素不可能都有这样唯一的幂级数展开。否则，人们就能像在函数域情形那样给出 $H^0(X_{S,E},\mathcal O_{Y_{S,E}}(n))$ 的简单描述。
\end{warning}

特别地，我们曾有 étale 层的短正合列
\[
  0 \to \bigcup_n G^{\mathrm{ad}}_{\breveE}[\pi^n] \to G^{\mathrm{ad}}_{\breveE} \xrightarrow{\log_G} \mathbb G_{a,\breveE} \to 0.
\]
过渡到万有覆盖，我们得到
\[
  0 \to V_\pi(G^{\mathrm{ad}}_{\breveE}) \to \tilde G^{\mathrm{ad}}_{\breveE} \xrightarrow{\log_G} \mathbb G_{a,\breveE} \to 0,
\]
其中 $V_\pi$ 是有理 $\pi$-进 Tate 模。由于 $G$ 的高度为 1，在几何点上有 $V_\pi(G^{\mathrm{ad}}_{\breveE}) \cong E$。作为 adic 空间，我们有
\[
  V_\pi(G^{\mathrm{ad}}_{\breveE}) \setminus \{0\} = \bigcup_{n \in \mathbb Z} \Spa \breveE_\infty .
\]

$V_\pi(G^{\mathrm{ad}}_{\breveE})$ 的一个元素涉及对所有 $n$ 相容地选取 $\pi^n$-根。特别地，给定一个 untilt $S^{\#}/\breveE_\infty$，我们得到一个特殊截面 $s \in \tilde G^{\mathrm{ad}}(S^{\#})$，它对应于所选相容的 $\pi^n$-挠点；这些挠点在对数下映到 0。于是在命题的同构下，我们得到一个映射
\[
  \mathcal O_{X_{S,E}} \to \mathcal O_{X_{S,E}}(1).
\]

\srcpage{36}它对应于 $s$。再与求值映射 $\mathcal O_{X_{S,E}}(1) \to \mathcal O_{S^{\#}}$ 复合，复合映射为 0，因为 $s$ 来自 $V_\pi \subset \ker(\log_G)$。所以上述映射穿过 $\mathcal O_{X_{S,E}} \to I_{S^{\#}}(1)$。

\begin{proposition}
这个映射 $\mathcal O_{X,S} \to I_{S^{\#}}(1)$ 是同构\yizhu{原文此处写作 $\mathcal O_{X,S}$，疑为 $\mathcal O_{X_{S,E}}$ 之误。}{A（原文 §7.7 p.36）}。
\end{proposition}

其实，只要以正确的方式思考，这就是形式论证。关键是理解上述等同（$V_\pi(G^{\mathrm{ad}}_{\breveE})\setminus\{0\} = \bigcup_{n\in\mathbb Z}\Spa\breveE_\infty$）。你归约到万有情形 $S^{\#} = \Spa \breveE_\infty$。此时 $Y_{S,E}$ 被典范地等同于 $\tilde G^{\mathrm{ad}}_{\breveE}\setminus\{0\}$，因为两者都是 $\Spa \mathcal O_{\breveE}[[\tilde X^{1/p^\infty}]] \setminus \{\pi = 0 \text{ 或 } \tilde x = 0\}$。然后你要分析截面 $\mathcal O_{X_{S,E}} \to \mathcal O_{X_{S,E}}(1)$ 的零点轨迹，并看出它在所有相关点处的阶为 1。之后你可以转到 $Y_{S,E}$ 上验证。

对于一般的 untilt（它只生活在 $\breveE$ 上），在 pro-étale 覆盖 $S^{\#}\times_{\breveE}\breveE_\infty \to S^{\#}$ 之后，我们得到等同 $I_{S^{\#}} \cong \mathcal O(-1)$。

这也完成了线丛的分类：untilt 的所有理想层都同构。

\begin{remark}
更一般地，我们可以考虑任何 $\pi$-可除 $\mathcal O_E$-模 $G/\ol{\FF}_q$。它带有 Dieudonné 模 $(V/\breveE,\phi_V)$。于是我们有 $\tilde G(S^{\#}) \cong \tilde G(S) \cong H^0(X_{S,E},\mathcal E(V))$。
\end{remark}
